\documentclass[10pt,a4paper]{article} 
\usepackage[utf8]{inputenc} 
\usepackage[T1]{fontenc} 
\usepackage[french]{babel} 
\usepackage{supertabular} %Nécessaire pour les longs tableaux
\usepackage[top=2.2cm, bottom=2.5cm, right=2.5cm, left=2.5cm]{geometry} %Mise en page 
\usepackage{amsmath} %Nécessaire pour les maths 
\usepackage{amssymb} %Nécessaire pour les maths 
%\usepackage{stmaryrd} %Utilisation des double crochets 
%\usepackage{pifont} %Utilisation des chiffres entourés 
\usepackage{graphicx} %Introduction d images 
\usepackage{epstopdf} %Utilisation des images .eps 
\usepackage{amsthm} %Nécessaire pour créer des théorèmes 
%\usepackage{algorithmic} %Nécessaire pour écrire des algorithmes 
%\usepackage{algorithm} %Idem 
\usepackage{bbold} %Nécessaire pour pouvoir écrire des indicatrices 
\usepackage{hyperref} %Nécessaire pour écrire des liens externes 
\usepackage{array} %Nécessaire pour faire des tableaux 
\usepackage{tabularx} %Nécessaire pour faire de longs tableaux 
\usepackage{caption} %Nécesaire pour mettre des titres aux tableaux (tabular) 
\usepackage{color} %nécessaire pour écrire en couleur 
\usepackage{float} % Pour l'option [H] de \begin{figure}
\newtheorem{thm}{Théorème} 
\newtheorem{mydef}{Définition}
\newtheorem{prop}{Proposition} 
\newtheorem{lemma}{Lemme}

\newcommand{\hmm}{\textsc{HMM}}
\newcommand{\mcmc}{\textsc{MCMC}}
\newcommand{\fhmm}{\textsc{Factorial HMM}}

\title{MVA, Projet PGM: Rapport initial\\
  Factorial HMM}
\author{Théis \textsc{Bazin} \and Valentin \textsc{De Bortoli} \and Élie \textsc{Michel}}

\begin{document}

\maketitle

Le but de ce rapport est de présenter l'avancement de nos travaux à la date du
9 décembre 2016.
Nous avons divisé notre présentation en trois parties : compréhension
mathématique du modèle, implémentation et choix des données.

\section{Modèle mathématique: Factorial HMM}

La plus grande partie de notre temps a été dédiée à la compréhension des
mathématiques sous-jacentes au modèle.
Nous commençons par une courte justification. Nous aurions pu modéliser des
données observables dépendant de $M$ variables cachées pouvant prendre $K$
valeurs comme un modèle probabiliste graphique de type \hmm,
\emph{Hidden Markov Model}, avec une variable cachée pouvant prendre
$K^M$ valeurs.

Néanmoins cela est très coûteux et ne prend pas en compte le découplage supposé
de nos $M$ variables cachées.
Pour cela les auteurs de l'article proposent une variante de \hmm :
\emph{Factorial HMM}.
Nous présentons en \autoref{fig:hmm} et \autoref{fig:factorial_hmm} les graphes
dirigés dans lesquels se factorisent les deux modèles.

\begin{figure}[hpbt]
  \centering
    \includegraphics[scale=0.6]{hmm.png}
  \caption{Hidden Markov Model\label{fig:hmm}}
\end{figure}

\begin{figure}[hpbt]
  \centering
    \includegraphics[scale=0.7]{facthmm.png}
  \caption{Factorial Hidden Markov Model\label{fig:factorial_hmm}}
\end{figure}

Il s'agit alors, comme dans le modèle \hmm{}, de faire de l'inférence sur
les paramètres de notre modèle.
Pour cela, nous utilisons l'algorithme EM.
L'étape de maximisation, \emph{M-step}, est bien comprise et très
similaire à celle de \hmm{}.
L'étape de calcul d'espérance, \emph{E-step}, est quant à elle différente
puisque nous ne pouvons a priori pas utiliser l'algorithme \emph{sum-product},
car il n'y a pas ici de structure d'arbre.

On note cependant que l'on a tout de même une structure particulière ($M$
couches de chaînes de Markov), ce qui permet d'établir un algorithme de calcul
exact des probabilités.
Malheureusement, ce calcul exact est très coûteux.
Nous considérons donc des approximations de l'espérance, basées sur trois
méthodes :

\begin{enumerate}
  \item L'échantillonnage de Gibbs,
  \item \emph{Mean-Field},
  \item \emph{Structured Mean-Field}.
\end{enumerate}

Les particularités mathématiques de chacun des modèles ont été étudiées.
L'échantillonnage de Gibbs est une implémentation classique d'un algorithme de
type \emph{Markov Chain Monte Carlo} (\mcmc).
\emph{Mean-Field} rend tous les états des variables cachées indépendants et
actualise les paramètres en conséquence.
Dans \emph{Structured Mean-Field}, on conserve la structure de chaîne de
Markov pour chaque couche, tout en retirant la dépendance entre chaque couche.

Un problème mathématique se pose : les paramètres du modèle sont actualisés
comme étant des points fixes d'une certaine fonction compliquée à calculer.
Les auteurs ne proposent pas de solution pour trouver ces points fixes, y
a-t-il une méthode exacte ou devons-nous utiliser des algorithmes itératifs ?

\section{Implémentation}

Concernant l'implémentation nous allons nous concentrer sur l'inférence de
paramètres du modèle.
Les algorithmes seront écrits en Python. L'implémentation n'a pas encore été
travaillée ; il nous a semblé plus judicieux de s'intéresser au devoir maison
qui traite du cas où $M=1$ : \hmm. La structure de notre code
s'organisera comme indiqué en \autoref{fig:implem_structure}.

\begin{figure}[hpbt]
  \centering
    \includegraphics[scale=1]{algograph.png}
  \caption{Organisation de l'implémentation\label{fig:implem_structure}}
\end{figure}

\section{Choix des données}

Tout d'abord, notre algorithme d'inférence de paramètres sera testé sur des
données simulées, puisqu'il est facile de faire un modèle génératif de
\fhmm.
Nous tenterons également de reproduire les résultats obtenus par les auteurs sur
les chorales de Bach.
Un autre but est également de tester notre algorithme sur des données réelles
issues de la sociologie (enquête 2013, "Génération 2010" sur l'évolution en
début de carrière de jeunes ayant quitté le système scolaire).
Les axes d'étude pour ce dernier jeu de données n'ont pas encore été établis.

\end{document}
